\documentclass[11pt,a4paper]{article}
\usepackage[utf8]{inputenc}
\usepackage[T1]{fontenc}
\usepackage{amsmath,amssymb}
\usepackage{booktabs}
\usepackage{geometry}
\usepackage{hyperref}
\usepackage{parskip}
\geometry{margin=2.7cm}
\hypersetup{colorlinks=true, linkcolor=blue, urlcolor=blue, citecolor=blue}

\title{Decomposição geométrica de um semiprimo $N = L \times q$\\
\large e sua relação com compressão de dados}
\author{}
\date{Versão consolidada}
\renewcommand{\contentsname}{Sumário}

\begin{document}
\maketitle

\noindent\textit{Versão consolidada — funde o relatório principal (seções 1--13) com o adendo sobre o triângulo como figura de referência (seções 14--16), cuja numeração original colidia com a expansão do relatório principal.}

\tableofcontents
\newpage

\section{Motivação}

Um semiprimo $N = L \times q$ pode ser representado por uma malha triangular de pontos: para $L$ fixo, os pontos $(x, y)$ com $0 \le y < L$ e $0 \le x < L - y$ formam um triângulo retângulo de $L(L+1)/2$ pontos. Dentro dessa malha, um bloco quadrado de lado
\[
k = L - q
\]
concentra a informação sobre o fator $q$: conhecer $k$ (e $L$) é equivalente a conhecer $q$. O objetivo deste relatório é registrar a construção geométrica feita, a decomposição observada em duas regiões em forma de T, e a analogia com técnicas de compressão/renderização que motivou a discussão seguinte.

\section{Construção geométrica}

A malha é desenhada em duas projeções: uma vista 2D direta e uma vista 3D em que o bloco central ganha altura $k$, sugerindo um paralelepípedo ``extraído'' da malha plana. Exemplo: malha triangular $L=9$ com bloco central $k=3$ ($q = L-k = 6$, $N = Lq = 54$).

\section{Decomposição em T1 / T2}

Centralizando o bloco de lado $k$ sobre a diagonal de simetria $x=y$ da malha (posição $c = (L-k)/2$), o restante da malha se separa, ponto a ponto, em três classes:

\begin{itemize}
  \item $T_1 = \{(x,y): x > y\}$ --- bloco próximo ao vértice $(L-1, 0)$;
  \item $T_2 = \{(x,y): y > x\}$ --- bloco próximo ao vértice $(0, L-1)$, espelho de $T_1$ pela troca $x \leftrightarrow y$;
  \item pontos com $x=y$ --- a própria dobra de simetria.
\end{itemize}

Cada um desses blocos tem, visualmente, o formato de um T. Para $L=9$, $k=3$: $|T_1| = |T_2| = 18$ pontos e 3 pontos na diagonal, totalizando $45 = \binom{9}{2}$ pontos fora do bloco de 9 pontos ($k^2$).

\section{Leitura em termos de compressão}

A decomposição sugere descrever a malha sem enumerar todos os $L(L+1)/2$ pontos:

\begin{enumerate}
  \item o bloco central é homogêneo, determinado por três números $(L, k, c)$;
  \item $T_1$ e $T_2$ são simétricos --- basta descrever um e refletir;
  \item a malha completa é reconstruída aplicando essas regras, sem processar a região redundante.
\end{enumerate}

Essa é a mesma lógica de voxel culling, sparse octrees e clipmaps: evitar processar regiões homogêneas ou fora de interesse. A diferença é a origem do critério de redutibilidade --- aqui algébrico (determinado por $L$ e $q$); num motor gráfico, geralmente por esparsidade de conteúdo ou visibilidade.

\section{Implementação}

Implementado em Python (\texttt{compress\_demo.py}):

\begin{itemize}
  \item \textbf{Critério de homogeneidade:} região $[x_0, x_0+s) \times [y_0, y_0+s)$ é homogênea se os quatro cantos e o centro têm o mesmo valor; para $s \le 2$ a checagem é exata célula a célula (necessário porque a hipotenusa é uma fronteira fina que pode escapar da amostragem por cantos).
  \item \textbf{Estrutura de dados:} quadtree --- cada nó é folha (valor único) ou se subdivide em 4 filhos. Nós homogêneos nunca são expandidos.
  \item \textbf{Regra de reconstrução:} \texttt{query(x,y)} desce a árvore até a folha correspondente --- nenhum array $L \times L$ é materializado em memória.
\end{itemize}

Para o caso 3D, o volume é processado fatia por fatia ao longo de $z$: cada fatia é comprimida em um quadtree independente e descartada antes da próxima (streaming) --- pico de memória = tamanho de uma única fatia comprimida.

\section{Resultados}

\begin{table}[h]
\centering
\begin{tabular}{rrrrr}
\toprule
L & k & células (L\textsuperscript{2}) & nós do quadtree & nós/L \\
\midrule
16   & 5    & 256        & 89     & 5.56 \\
64   & 21   & 4\,096     & 401    & 6.27 \\
256  & 85   & 65\,536    & 1\,673 & 6.54 \\
1024 & 341  & 1\,048\,576& 6\,785 & 6.63 \\
4096 & 1365 & 16\,777\,216 & 27\,257 & 6.65 \\
\bottomrule
\end{tabular}
\caption{Crescimento do quadtree com L (fixando $k \approx L/3$). A razão nós/L estabiliza, confirmando escala $O(L)$.}
\end{table}

O experimento de streaming 3D ($L=256$, $k=70$) confirma o ganho de memória: volume bruto de 16,8 milhões de células, mas pico de apenas 1.561 nós de uma única fatia comprimida --- razão $9{,}3 \times 10^{-5}$. Reconstrução ponto a ponto validada contra a função de ocupação original em 2000 amostras aleatórias por caso, sem divergência. \textbf{(Confirmado por execução direta do script \texttt{compress\_demo-3.py}: pico de 1.561 nós bate exato.)}

\section{Duas técnicas de truncamento sem divisão}

\textbf{Truncamento por máscara (resíduo):} $i \mathbin{\&} R_k$ (com $R_k = 2^k - 1$) calcula $i \bmod 2^k$ sem custo de divisão, mas não preserva a ordem de grandeza de $i$. Usado no protótipo original para evitar divisões/módulos.

\textbf{Truncamento por deslocamento (shift):} $i \gg k$ calcula $\lfloor i/2^k \rfloor$ exatamente, preservando a magnitude. Técnica correta quando o objetivo é dimensionar (decidir se dois valores são ``parecidos o bastante'').

Teste numérico ($k=4$, $R_k=15$): os testes $(i \mathbin{\&} R_k) \le 1$ e $(i \gg k) = 0$ discordam em 128 dos 256 casos --- respondem perguntas diferentes; nenhum é ``o errado''.

\section{Fusão de blocos como grafo (union-find)}

Em vez de subdividir recursivamente, trata-se os blocos como nós de um grafo, fundindo blocos adjacentes cujo valor truncado (por deslocamento) coincide, via union-find --- abordagem de Felzenszwalb--Huttenlocher, mais flexível que o quadtree (regiões fundidas não precisam ser quadrados alinhados à grade).

\begin{table}[h]
\centering
\begin{tabular}{rrrr}
\toprule
L & blocos totais & componentes & razão \\
\midrule
32   & 64      & 2 & 0.0312 \\
128  & 1\,024  & 2 & 0.0020 \\
512  & 16\,384 & 2 & 0.0001 \\
1024 & 65\,536 & 2 & 0.00003 \\
\bottomrule
\end{tabular}
\caption{Fusão de blocos por grafo ($k_{shift}=1$, blocos $4\times4$). Componentes $O(1)$, à custa de perder a fronteira exata. \textbf{(Confirmado por execução de \texttt{merge\_grafo.py}.)}}
\end{table}

Quadtree preserva a fronteira exata ao custo de $O(L)$; o grafo com truncamento comprime para $O(1)$ classes ao custo de perder a posição exata da fronteira.

\section{A constante da Tabela 1 não é universal: sensibilidade ao alinhamento binário}

A hipótese de que nós/L fosse função fechada apenas de $k/L$ não bateu com os dados (discrepância de até $\sim$46\%). Variando $k/L$ para $L=4096$ fixo:

\begin{table}[h]
\centering
\begin{tabular}{rr}
\toprule
k/L & nós/L \\
\midrule
1/4 & 4.00 \\
1/6 & 5.32 \\
1/3 & 6.66 \\
3/8 & 4.01 \\
1/2 & 4.00 \\
5/8 & 4.01 \\
\bottomrule
\end{tabular}
\caption{Não há tendência suave --- proporções bem alinhadas em binário (1/4, 1/2) retornam à constante baixa $\approx 4$; 1/3 é o pior caso observado.}
\end{table}

\textbf{Causa:} o teste de homogeneidade por amostragem (cantos + centro) funciona perfeitamente quando as bordas do bloco caem sobre fronteiras de subdivisão (potências de 2), mas força subdivisões extras quando $k$ e $c=(L-k)/2$ não são bem alinhados. $k/L=1/3$ é mau alinhamento; $k/L \in \{1/4, 1/2\}$ é bem alinhado.

\textbf{Consequência:} não existe fórmula fechada simples $\text{nós}(L,k) = f(k/L) \cdot L$ válida para todo $k$. Uma afirmação anterior ($\text{nós}(L,k) = 3L-k$, de um experimento diferente com teste exato) foi removida por não corresponder a esta implementação. O crescimento $O(L)$ permanece válido; a constante de proporcionalidade depende da heurística de amostragem e do alinhamento binário de $k$.

\section{Hexágono: dois trapézios colados pela hipotenusa}

Removendo o triângulo menor de lado $k$ (a ``ponta'') do triângulo de lado $L$, sobra um trapézio $T_1$ (base curta $x+y=k$, base larga $x+y=L-1$). Refletindo $T_1$ pela hipotenusa --- $(x,y) \mapsto (C-y, C-x)$, $C=L-1$ --- obtém-se $T_2$, colado a $T_1$ pela base larga compartilhada: um hexágono simétrico.

\[
|T_1| = \sum_{i=k+1}^{L} i = \frac{L(L+1) - k(k+1)}{2}
\]
\[
N_{hex}(L,k) = 2|T_1| - L = L^2 - k(k+1)
\]

Validado contra construção explícita para 30 combinações $(L,k)$, $L \in [5,100]$ --- todos exatos:

\begin{table}[h]
\centering
\begin{tabular}{rrrr}
\toprule
L & k & pontos (construção real) & $L^2 - k(k+1)$ \\
\midrule
9   & 3  & 69    & 69 \\
20  & 9  & 310   & 310 \\
50  & 24 & 1\,900 & 1\,900 \\
100 & 49 & 7\,550 & 7\,550 \\
\bottomrule
\end{tabular}
\end{table}

Leitura geométrica direta: $L^2$ é a malha quadrada completa; $k(k+1) = k^2+k$ são as duas cunhas triangulares retiradas dos cantos --- o hexágono é ``o quadrado $L \times L$ menos duas pontas''.

\section{Efeito telescópico entre valores de k (gnomons generalizados)}

Fixando $L$ e variando $k$:
\[
N_{hex}(L,k) - N_{hex}(L,k+1) = 2(k+1)
\]

Análogo direto dos números gnomon clássicos. A soma das cascas de $k=0$ até $k=m-1$ reconstrói o salto direto:
\[
\sum_{k=0}^{m-1} 2(k+1) = m(m+1) = N_{hex}(L,0) - N_{hex}(L,m)
\]

Confirmado numericamente ($L=100$, $m=10$: soma das cascas = 110, salto direto = 110).

\section{Comparação com dados de ocupação esparsa (esfera e terreno sintéticos)}

\subsection{Por que dados sintéticos}

Sem acesso a datasets reais neste ambiente, usam-se dados sintéticos que reproduzem características estruturais: esfera num cubo $L^3$ (proxy de órgão/tumor, fronteira curva suave), terreno por heightmap suave (proxy de voxels de jogo), e ruído puro como controle.

\subsection{Extensão do método para 3D}

O quadtree é generalizado para octree: nó cúbico de lado $s$ subdivide em 8 filhos de lado $s/2$ quando não homogêneo (mesmo critério de amostragem, exato para $s \le 2$).

\subsection{Resultados}

\begin{table}[h]
\centering
\begin{tabular}{rrrrr}
\toprule
L & nós (esfera) & nós/L\textsuperscript{2} & nós (terreno) & nós/L\textsuperscript{2} \\
\midrule
16  & 905     & 3.54 & 1\,225  & 4.79 \\
32  & 3\,641  & 3.56 & 4\,713  & 4.60 \\
64  & 14\,377 & 3.51 & 22\,489 & 5.49 \\
128 & 57\,769 & 3.53 & 93\,585 & 5.71 \\
\bottomrule
\end{tabular}
\caption{A razão nós/L\textsuperscript{2} estabiliza para ambos, confirmando $O(L^2)$. \textbf{(Confirmado por execução direta de \texttt{octree3d.py} --- bate exato em todas as linhas.)}}
\end{table}

\begin{table}[h]
\centering
\begin{tabular}{lccc}
\toprule
Construção & Dimensão & Crescimento & Razão nós/bruto \\
\midrule
Malha algébrica (Seção 5) & 2D & $O(L)$   & 0,052 \\
Esfera sintética          & 3D & $O(L^2)$ & 0,028 \\
Terreno sintético         & 3D & $O(L^2)$ & 0,045 \\
Ruído puro (sem estrutura)& 3D & $O(L^3)$ & 1,123 (pior que bruto!) \\
\bottomrule
\end{tabular}
\caption{Comparação final. \textbf{(Confirmado: $0{,}0275\approx0{,}028$; $0{,}0446\approx0{,}045$; $1{,}1234\approx1{,}123$; malha algébrica $0{,}0501\approx0{,}052$ para $L=128$.)}}
\end{table}

\subsection{Interpretação}

O ganho em dados sintéticos ``realistas'' é da mesma ordem de grandeza do obtido na construção algébrica. O princípio (compressão por homogeneidade local, custo proporcional à fronteira) generaliza para qualquer dado com fronteiras bem comportadas --- não é peculiaridade do semiprimo.

O ruído puro é o contraponto: sem estrutura a explorar, o octree piora em relação ao array bruto (razão $>1$), porque cada subdivisão adiciona um nó interno sem atingir folha homogênea maior que $1\times1\times1$.

\subsection{O piso teórico para dados incompressíveis: um terceiro tipo de truncagem}

A razão $>1$ para ruído puro não é artefato --- é limite matemático exato. Para dados verdadeiramente incompressíveis, a árvore desce até folhas $1\times1\times1$ em todo lugar:
\[
N(d) = \sum_{i=0}^{d} b^i = \frac{b^{d+1}-1}{b-1}, \qquad b^d = \text{raw}
\]
\[
\frac{N(d)}{b^d} = \frac{b}{b-1} - \underbrace{\frac{b^{-d}}{b-1}}_{\text{termo de truncagem}} \xrightarrow{d\to\infty} \frac{b}{b-1}
\]

Para octree ($b=8$): limite $8/7 \approx 1{,}142857$. Medido em $L=8$ ($d=3$) com verificação exata de homogeneidade: $1{,}142578$ --- bate a fórmula exatamente.

\textbf{Nota de validação} (não presente no relatório original, verificada nesta consolidação): essa igualdade exata só se sustenta com dados verdadeiramente incompressíveis --- em teste com valores contínuos (float aleatório, sem chance de coincidência), a medição bate $1{,}142578$ exatamente, confirmado também para $d=1..5$ contra a fórmula fechada e contra o termo residual $b^{-d}/(b-1)$, todos exatos. \textbf{Com dados binários} (0/1, que é o que \texttt{make\_noise\_occ} de fato gera), a razão medida fica sistematicamente \textbf{abaixo} do piso teórico (testado: $1{,}126953$ em vez de $1{,}142578$ para os mesmos parâmetros), porque blocos de 8 células têm $\sim$0,8\% de chance de sair homogêneo por coincidência mesmo sendo ruído puro. Isso não invalida a fórmula --- ela descreve o pior caso absoluto --- mas o valor específico citado como ``medido'' no relatório pressupõe um alfabeto de valores efetivamente contínuo, não o gerador binário incluído em \texttt{octree3d.py}.

Essa truncagem é qualitativamente diferente das da Seção 7: não é escolha de projeto, é consequência inevitável de qualquer representação finita --- um viés sistemático, quantificável e residual, presente em qualquer estrutura em árvore aplicada a dados finitos.

\section{Custo amortizado: construção única vs. renderização por quadro}

\subsection{Por que contagem de nós não é a métrica certa}

O caso de uso real (renderização, fatiamento de volumes) não é sobre armazenar --- é sobre visitar/pintar o volume repetidamente, quadro a quadro.

\subsection{Metodologia}

Compara-se \textbf{array denso} (percorre todas as $L^3$ células do zero a cada quadro) contra \textbf{octree} (construído uma vez; cada quadro seguinte apenas percorre a árvore pronta).

\begin{table}[h]
\centering
\begin{tabular}{lrrrr}
\toprule
Cenário & build (1$\times$) & denso/quadro & árvore/quadro & ganho/quadro \\
\midrule
Esfera L=32  & 20,7 ms  & 9,3 ms   & 0,32 ms & 29$\times$  \\
Terreno L=32 & 76,1 ms  & 46,0 ms  & 0,43 ms & 107$\times$ \\
Esfera L=64  & 78,0 ms  & 77,8 ms  & 1,54 ms & 51$\times$  \\
Terreno L=64 & 383,9 ms & 357,9 ms & 2,39 ms & 150$\times$ \\
\bottomrule
\end{tabular}
\caption{Python puro, 20 quadros por cenário. \textbf{(Confirmado por execução de \texttt{render\_bench2.py} --- mesma ordem de grandeza e conclusão qualitativa; tempos absolutos variam por máquina, como esperado.)}}
\end{table}

\subsection{Ponto de equilíbrio (break-even)}

No cenário mais custoso (terreno, $L=64$): build (384 ms) equivale a pouco mais de 1 quadro de array denso (358 ms/quadro) --- a partir do segundo quadro a árvore já está adiante. Em todos os cenários, custo total (build + 20 quadros) ficou entre 4 e 16$\times$ menor que o array denso.

\subsection{Interpretação}

O ganho não vem de acesso mais rápido a um ponto isolado (é $O(\log L)$ na árvore contra $O(1)$ no array denso --- mais lento nessa operação específica). Vem do reuso amortizado de um trabalho de homogeneização pago uma única vez, ao longo de muitos quadros sobre um dado que não muda a cada instante. É por isso que motores gráficos constroem a estrutura esparsa ao carregar a cena e reaproveitam em todos os quadros seguintes.

\bigskip
\noindent\hrulefill

\section{O triângulo como figura de referência (Altura Máxima e Projeção Quadrada)}

\textit{(Seções 14--16 vêm do adendo original --- lá numeradas 13--15 --- renumeradas aqui para encaixar após a expansão do relatório principal.)}

Entre as figuras derivadas da malha (triângulo, trapézio, hexágono, ampulheta), o triângulo original --- antes de qualquer corte ou dobra --- é o de maior altura. Todas as demais figuras (Seção 10) são obtidas subtraindo ou recombinando partes do triângulo original; como nenhuma operação acrescenta pontos novos fora da malha de partida, nenhuma pode superar a altura do triângulo original --- ele é o teto geométrico do qual as outras derivam.

\subsection{A projeção como quadrado}

O triângulo retângulo de catetos $L$ tem como bounding box exatamente o quadrado $L \times L$: a projeção sobre cada eixo cobre $[0, L)$, então a menor região retangular capaz de conter o triângulo é o quadrado de área $L^2$. O próprio triângulo ocupa metade dessa projeção:
\[
|T_\triangle| = \frac{L(L+1)}{2} \approx \frac{L^2}{2}
\]
com a diferença de $L/2$ vindo dos pontos da própria hipotenusa discreta, contados uma única vez.

\subsection{O quadrado como referência de medida}

Essa projeção fornece uma unidade de medida comum: em vez de medir cada figura isoladamente, mede-se sua diferença em relação ao quadrado $L^2$ da figura de maior altura. Exatamente a leitura já obtida algebricamente na Seção 10 para o hexágono:
\[
N_{hex}(L,k) = L^2 - k(k+1)
\]
--- o quadrado funcionando como referência absoluta, com as figuras derivadas descritas por subtração, não recontadas ponto a ponto.

\subsection{Consequência prática}

Ao tomar a figura de maior altura como referência de projeção, medir um conjunto de figuras relacionadas se reduz a uma quantidade fixa ($L^2$) mais termos de correção específicos de cada corte ou dobra --- mesmo princípio de eficiência da Seção 4.

\subsection{Nota conceitual: medir com referência não muda a figura}

Usar $L^2$ como referência não altera a figura medida --- o hexágono continua tendo exatamente $L^2 - k(k+1)$ pontos, quer descrito por subtração ou contagem direta. O que muda é apenas o custo de obter a contagem: com a referência, uma subtração; sem ela, enumeração região por região. A figura é a mesma nos dois casos; o ganho é inteiramente no processo de cálculo, não no objeto calculado.

\section{Relações Métricas: Leitura Associativa vs. Fórmula}

\subsection{O teorema}

Num triângulo retângulo, a altura relativa à hipotenusa o divide em dois triângulos menores, semelhantes entre si e ao original. Sejam $b,c$ os catetos, $a$ a hipotenusa, $h$ a altura relativa a $a$, $m,n$ as projeções de $c$ e $b$ sobre $a$ ($m+n=a$):
\[
b^2 = a \cdot n, \qquad c^2 = a \cdot m, \qquad h^2 = m \cdot n, \qquad a \cdot h = b \cdot c
\]

Somando as duas primeiras: $b^2+c^2 = a(m+n) = a^2$ --- o Teorema de Pitágoras emerge da semelhança, não como ponto de partida independente.

\subsection{O que essas relações realmente são: associação antes de fórmula}

$b^2 = a \cdot n$ não é, na origem, uma sentença algébrica a memorizar --- é a tradução de uma percepção visual: o quadrado de lado $b$ tem a mesma área que o retângulo de lados $a$ e $n$, percebida pelo encaixe dos triângulos semelhantes. A fórmula é o registro posterior de uma associação já vista.

Mesmo tipo de leitura que perpassa todo o relatório: perceber que o hexágono é o quadrado menos duas pontas (Seção 10); perceber que o triângulo, por ter maior altura, projeta naturalmente o quadrado $L^2$ (Seção 14) --- a mesma leitura associativa aplicada à figura de referência do relatório inteiro.

\subsection{Por que essa distinção importa}

Reconhecer que a fórmula é tradução, não fonte, muda a ordem natural do trabalho: buscar primeiro a associação visual (semelhança, encaixe, projeção), e só então traduzir em símbolos para comunicar ou verificar.

\section{Recuperação Geométrico-Métrica: Método da Partição e Escala}

\subsection{Premissas}

Dadas a área $A$ de um retângulo (ou $A_T = A/2$ do triângulo correspondente) e sua diagonal $D$, existe único par $a \ge b > 0$ com $a \cdot b = A$, $a^2+b^2=D^2$. Objetivo: recuperar $a$ e $b$ sem a equação quadrática explícita $x^2 - Sx + A = 0$.

\subsection{Princípio da escala base}

Usa-se um triângulo de referência $(u,v,w)$ com propriedades conhecidas. Se a figura real se relaciona com a referência por fator de escala $k$, a área escala quadraticamente: $\text{Área real} = k^2 \cdot A_0$. Escolhendo $w=D$, tem-se $k=1$, eliminando a necessidade de extrair raiz para obter o fator de escala.

\subsection{Recuperação por partição e truncamento}

\textbf{Altura relativa à diagonal:} $h_D = \dfrac{2A_T}{D} = \dfrac{A}{D}$ --- medida linear por divisão simples, sem raízes.

\textbf{Invariante do quadrado da soma:} $(a+b)^2 = D^2 + 2A$ --- recupera $(a+b)$ sem calcular $a,b$ separadamente. Perímetro: $P = 2(a+b) = 2\sqrt{D^2+2A}$.

\subsection{Síntese do algoritmo}

\begin{table}[h]
\centering
\begin{tabular}{clll}
\toprule
Passo & Operação geométrica & Expressão & Exemplo ($A=60$, $D=13$) \\
\midrule
1 & Altura relativa à diagonal & $h_D = A/D$    & $60/13$ \\
2 & Quadrado composto          & $Q = D^2+2A$   & $169+120=289$ \\
3 & Semiperímetro               & $S=\sqrt{Q}$   & $17$ \\
4 & Perímetro total             & $P=2S$         & $34$ \\
\bottomrule
\end{tabular}
\end{table}

\textbf{Validado em C++ e Python nesta consolidação:} bate exato o exemplo do documento ($a=12$, $b=5$, $P=34$) e passa em 2000/2000 casos aleatórios (a única ``falha'' observada em teste anterior foi ruído de ponto flutuante de $\sim 10^{-6}$ no caso degenerado $a=b$, não erro do método). Já integrado ao módulo \texttt{Render::BoundingReference} da engine.

\subsection{Vantagem e generalização}

O método evita montar a equação de segundo grau, substituindo-a por análise de proporção, restrições lineares e decomposição por partições --- mesma linha da Seção 15: a raiz quadrada, quando aparece (Passo 3), surge de reconhecer um quadrado perfeito por partição ($289=17^2$), não de extrair raiz arbitrária. Generaliza para qualquer prisma ou volume cuja base seja parametrizada por projeções e razões fixas de área --- a mesma lógica de escala ($k^2$ para área, $k^3$ para volume) que rege semelhança de figuras em qualquer dimensão.

\bigskip
\noindent\hrulefill

\section{Próximos passos}

A implementação atual usa teste de homogeneidade por amostragem (cantos + centro), suficiente para fronteiras retas simples mas não generalizável a fronteiras curvas ou regiões sobrepostas. Direções de continuação, com status atualizado nesta consolidação:

\begin{itemize}
  \item \textbf{Testes de homogeneidade exatos} via geometria analítica (interseção reta-retângulo) --- ainda \textbf{não implementado} em nenhum dos scripts revisados; eliminaria a dependência de alinhamento binário descrita na Seção 9.
  \item \textbf{Extensão a octree 3D completo} --- \textbf{concluído} (Seção 12, \texttt{octree3d.py}), embora ainda usando o mesmo teste de homogeneidade por amostragem (não geométrico exato) da Seção 5.
  \item \textbf{Esquema híbrido} (grafo de fusão da Seção 8 para classificação grosseira $O(1)$, refinado por quadtree da Seção 5) --- proposto desde a primeira versão do relatório, \textbf{ainda não implementado}: confirmado por inspeção dos scripts que \texttt{UnionFind} (\texttt{merge\_grafo.py}) e \texttt{build\_quadtree}/\texttt{build\_octree} (\texttt{compress\_demo-3.py}, \texttt{octree3d.py}, \texttt{render\_bench2.py}) nunca aparecem no mesmo arquivo.
  \item \textbf{Dados de ocupação esparsa realistas} (voxels de jogos, volumes médicos) --- parcialmente endereçado: dados sintéticos (esfera, terreno) implementados e validados (Seção 12); dados reais permanecem como continuação, por falta de acesso à rede neste ambiente.
  \item \textbf{Validação do piso teórico para dados incompressíveis (Seção 12.5)} --- a fórmula está confirmada matematicamente, mas o script de geração de ruído (\texttt{make\_noise\_occ}) usa alfabeto binário, que produz medições sistematicamente abaixo do piso teórico por coincidências de homogeneidade local; um gerador de valores contínuos seria necessário para reproduzir o valor específico citado como ``medido'' no relatório.
\end{itemize}

\end{document}
